% example-traffic-light.tex  --  Section 5.1, traffic-light example.
% Follows Examples/Models/TrafficLight.lean (label, controller, stuckController,
% safety, eventuallyGreen, redResponds, stuckPath). Nodes are controller states;
% {red}, {green}, {yellow} are proposition labels and not automaton states.
%   (a) normal cycle   
%   (b) the same cycle plus the single extra edge stop -> stop

\begin{tikzpicture}[x=1mm, y=1mm,
  tlstate/.style={kiltstage, circle, draw=kiltnavy, fill=kiltpale, line width=0.8pt,
    minimum size=13mm, inner sep=0pt, font=\footnotesize\sffamily, text=kiltnavy},
  tlprop/.style={font=\footnotesize, text=kiltblue, inner sep=1pt, anchor=west},
  tlblock/.style={font=\footnotesize, text=kiltnavy, inner sep=0pt, align=left,
    anchor=north west, text width=39mm},
  tltitle/.style={font=\footnotesize\bfseries, text=kiltnavy, anchor=west, inner sep=0pt},
]

% =====================================================================
% (a)  normal labelled cycle                           x = 0 .. 76 mm
% =====================================================================
\node[tltitle] at (0,41) {(a) Normal labelled cycle};

\node[tlstate] (s) at (12.5,18)  {stop};
\node[tlstate] (g) at (30.5,30.5) {go};
\node[tlstate] (c) at (30.5,6.5)  {caution};

\draw[kiltarrow] (0.3,18) -- (s.west);          % initial
\draw[kiltarrow] (s) -- (g);
\draw[kiltarrow] (g) -- (c);
\draw[kiltarrow] (c) -- (s);

\node[tlprop, anchor=north] at (12.5,10.5) {$\{\mathit{red}\}$};
\node[tlprop] at (37.5,30.5) {$\{\mathit{green}\}$};
\node[tlprop] at (37.5,6.5)  {$\{\mathit{yellow}\}$};

\node[tlblock] at (38,26.5)
  {exactly one colour\\
   eventual green holds\\
   red response holds};

% =====================================================================
% (b)  the added self-loop                              x = 82 .. 158 mm
% =====================================================================
\begin{scope}[shift={(82,0)}]
\node[tltitle] at (0,41) {(b) One added self-loop};

\node[tlstate, draw=kiltorange, fill=kiltpeach, text=kiltorange] (s) at (12.5,18) {stop};
\node[tlstate] (g) at (30.5,30.5) {go};
\node[tlstate] (c) at (30.5,6.5)  {caution};

\draw[kiltarrow] (0.3,18) -- (s.west);          % initial
\draw[kiltarrow, kiltblue!40] (s) -- (g);        % unchosen edge: faded, still visible
\draw[kiltarrow] (g) -- (c);
\draw[kiltarrow] (c) -- (s);
\draw[kiltarrow, kiltdashed, kiltorange, line width=1.3pt]
  (s.150) to[out=150,in=85,looseness=4.5] (s.85);                      % stop -> stop

\node[font=\footnotesize, text=kiltorange, inner sep=0pt] at (7,36) {added loop};

\node[tlprop, anchor=north] at (12.5,10.5) {$\{\mathit{red}\}$};
\node[tlprop] at (37.5,30.5) {$\{\mathit{green}\}$};
\node[tlprop] at (37.5,6.5)  {$\{\mathit{yellow}\}$};

\node[tlblock] at (38,26.5)
  {$\pi=\textsf{stop}^{\omega}$\\
   one-colour safety holds\\
   \textcolor{kiltorange}{eventual green fails}\\
   \textcolor{kiltorange}{red response fails}};
\end{scope}

\node[font=\footnotesize, text=kiltblue, inner sep=0pt] at (79,-4)
  {Both variants are serial; no fairness is assumed.};

\end{tikzpicture}
